\section{Research Problem and Gap}

% ---------------------------------------------------------------------------------------------------------------------
\subsection{Limitations of existing moderation tools in remote settings}
Current tools fall short in three key areas:
\begin{itemize}
  \item Passive functionality: Most AI tools (e.g., Zoom AI Companion, Google Meet’s recaps) focus on
    post-meeting summaries or transcription, not real-time moderation [Zoom, 2024; Google, 2023].
  \item Generic interventions: Solutions like Miro’s AI facilitator provide template-based guidance but fail
    to adapt to the meeting’s context or user preferences [Miro, 2024].
  \item Lack of empirical validation: Few studies evaluate user acceptance of proactive AI moderation,
    particularly in high-stakes decision-making scenarios [Benbasat and Barki, 2007].
\end{itemize}

% ---------------------------------------------------------------------------------------------------------------------
\subsection{Need for context-sensitive, proactive AI systems }
To address these gaps, this thesis proposes a proactive AI moderation system that:
\begin{itemize}
  \item Detects thematic deviations in real time using NLP and meeting goal tracking.
  \item Intervenes constructively with personalized, non-intrusive prompts (e.g., private messages to
    individuals or group reminders).
  \item Validates design choices through user studies focusing on usability (SUS) and acceptance (TAM).
\end{itemize}
Such a system could reduce meeting fatigue, improve decision quality, and lower the barrier to effective
facilitation—but its success hinges on user trust and perceived usefulness, which remain underexplored in literature.
